
\BookRecordIndex{T0003}$J$-全纯曲线理论由\personfirst{geluomofu}{格罗莫夫}{Gromov}于 1985 年创立，
它深刻影响了辛几何的研究，如今几乎渗透到辛几何的各个方面。这些方法对\BookRecordIndex{T0207}凯勒流形的
研究也很有意义，因为在研究这类流形中全纯曲线的性质时，常常适宜把复结构扰动成一般
位置。这样做可保证我们看到的是这些曲线稳定存在的特征，而不是偶然出现的特征。不过，
扰动后的结构可能不再可积，于是我们仍会进入这样一类曲线的研究：它们相对于某个不可积
\BookRecordIndex{T0088}殆复结构 $J$ 是全纯的。

这些曲线满足\BookRecordIndex{T0209}柯西--黎曼方程的一种非线性类似物。在能够用它们做任何有用的事情之前，
必须先掌握这些方程理论的基本内容；例如，必须知道哪些条件能保证解空间是有限维流形，
还要知道如何处理这些解空间通常并不紧致这一事实。如前言所述，本章介绍全部基本概念，
并概述建立上述结果所需的要素。读者，尤其是不熟悉这一理论的读者，应当先读本章，随后
再分别深入研读自己感兴趣的方面。

我们假定读者熟悉辛几何的基础。已有若干入门书籍可供参考。一种选择是
McDuff--Salamon \srcite{277}，此外现在还有 Berndt \srcite{35}、
Cannas da Silva \srcite{59} 等更初等的论述，以及 Arnold \srcite{15} 等经典著作。

\section{\BookRecordIndex{T0358}辛流形}\label{sec:1-1}
\zhindex{xinliuxing}{辛流形}

光滑 $2n$ 维流形 $M$ 上的一个辛结构，是一个闭 $2$-形式 $\omega$，它是非退化的，
这里“非退化”是指最高维形式 $\omega^n$ 处处不消失。由达布（Darboux）定理，所有
\BookRecordIndex{T0363}辛形式在局部都微分同胚于欧氏空间 $\mathbb R^{2n}$ 上的标准\BookRecordIndex{T0363}辛形式
\[
  \omega_0=dx_1\wedge dy_1+\dots+dx_n\wedge dy_n,
\]
其中坐标为 $(x_1,y_1,\dots,x_n,y_n)$。这使得\BookRecordIndex{T0358}辛流形的整体结构难以把握。人们发展了
变分技术，用来研究欧氏空间以及余切丛等具有某种线性结构的流形中的若干整体问题；参见
\srcite{190, 277} 及其中所列参考文献。不过，在最广泛的一类\BookRecordIndex{T0358}辛流形上都适用的方法，
是 \BookRecordIndex{T0003}$J$-全纯曲线方法。

这里 $J$ 是 $M$ 上一个被 $\omega$ \BookRecordIndex{T0372}驯服的\BookRecordIndex{T0088}殆复结构。\BookRecordIndex{T0088}殆复结构是 $M$ 的切丛 $TM$
的一个自同构 $J$，满足恒等式
\[
  J^2=-\mathbf 1.
\]
因此可以把 $J$ 看成乘以 $i$，它使 $TM$ 成为一个复维数为 $n$ 的复向量丛。若对一切
非零 $v\in TM$ 都有
\[
  \omega(v,Jv)>0,
\]
就称 $\omega$ \BookRecordIndex{T0372}驯服 $J$。在几何上，这表示 $\omega$ 在切空间 $T_xM$ 的每条复直线
$L=\operatorname{span}\{v,Jv\}$ 上的限制都是正形式。给定 $\omega$，所有被 $\omega$
\BookRecordIndex{T0372}驯服的\BookRecordIndex{T0088}殆复结构组成的集合 $\mathcal J_\tau(M,\omega)$ 总是非空且可缩。请注意，构造和
扰动被\BookRecordIndex{T0372}驯服的\BookRecordIndex{T0088}殆复结构非常容易，因为定义它们的条件是逐点的。还要注意，由于
$\mathcal J_\tau(M,\omega)$ 道路连通，不同的 $J\in\mathcal J_\tau(M,\omega)$ 给出
彼此同构的复向量丛 $(TM,J)$。因此这些丛的陈类与 $J$ 的选择无关，记为 $c_i(TM)$。
\footnote{与 $(M,\omega)$ 自然相伴的还有另一个\BookRecordIndex{T0088}殆复结构空间，即第 3.1 节定义的
$\omega$-\BookRecordIndex{T0348}相容结构空间 $\mathcal J(M,\omega)$。就当前目的而言，两者均可使用。事实上，
为了使结果尽可能一般，我们常在 $\mathcal J(M,\omega)$ 上工作，因为这略难一点：
$\mathcal J_\tau(M,\omega)$ 在所有\BookRecordIndex{T0088}殆复结构组成的空间中是开集，而
$\mathcal J(M,\omega)$ 不是。}

下文只需用到第一陈类；有关的是它在同调类 $A\in H_2(M)$ 上所取的值
\[
  c_1(A):=\langle c_1(TM),A\rangle.
\]
若 $A$ 由定义在闭定向 $2$-流形 $\Sigma$ 上的光滑映射 $u:\Sigma\to M$ 表示，
那么 $c_1(A)=c_1(E)$ 是拉回切丛 $E:=u^*TM\to\Sigma$ 的第一陈数。可是，$2$-流形
$\Sigma$ 上的每个复向量丛 $E$ 都\BookRecordIndex{T0213}可分解为复线丛之和
$E=L_1\oplus\dots\oplus L_n$。相应地，
\[
  c_1(E)=\sum_i c_1(L_i).
\]
复线丛的第一陈数恰好就是它的欧拉数，因此往往很容易直接计算 $c_1(L_i)$。例如，若
$A$ 是 $M=V\times S^2$ 中球面 $S=\mathrm{pt}\times S^2$ 的类，那么显然
\[
  TM\big|_S=TS\oplus L_2\oplus\dots\oplus L_n,
\]
其中线丛 $L_k$ 都是平凡的。于是
\[
  c_1(A)=c_1\bigl(TM\big|_S\bigr)=c_1(TS)=\chi(S)=2,
\]
这里 $\chi(S)$ 是 $S$ 的欧拉示性数。

从一个殆复流形到另一个殆复流形的光滑映射
$\phi:(M,J)\to(M',J')$，称为 $(J,J')$-全纯的，当且仅当其微分
$d\phi(x):T_xM\to T_{\phi(x)}M'$ 是复线性的，即
\[
  d\phi(x)\circ J(x)=J'(\phi(x))\circ d\phi(x).
\]
这些就是\BookRecordIndex{T0209}柯西--黎曼方程。当 $(M,J)$ 和 $(M',J')$ 都是复 $n$ 维空间
$\mathbb C^n$ 的子集时，满足这些方程的恰好是全纯映射。如果\BookRecordIndex{T0088}殆复结构 $J$ 来自 $M$
上的一个底\BookRecordIndex{T0072}层复结构，就称 $J$ 可积。等价地说，可以为 $M$ 选择一个图册，使其坐标卡
都是 $(J,i)$-全纯的，其中 $i$ 是 $\mathbb C^n$ 上的标准复结构。此时坐标变换是
$\mathbb C^n$ 中开集之间通常意义下的全纯映射。当 $M$ 为 $2$ 维时，一个基本定理断言
$M$ 上每个\BookRecordIndex{T0088}殆复结构 $J$ 都可积；证明见第 E.4 节。但在更高维数中，情况远非如此。

欧氏空间 $(\mathbb R^{2n},\omega_0)$ 连同其标准\BookRecordIndex{T0088}殆复结构 $J_0$，是殆复\BookRecordIndex{T0358}辛流形的
基本例子。这里 $J_0$ 来自通常的同一化 $\mathbb R^{2n}=\mathbb C^n$：
\[
  z_j=x_j+iy_j,\qquad j=1,\dots,n.
\]
因而在 $\mathbb R^{2n}=T_z\mathbb R^{2n}$ 的标准基下，$J_0$ 把切向量
$\partial/\partial x_j$ 映到 $\partial/\partial y_j$，把
$\partial/\partial y_j$ 映到 $-\partial/\partial x_j$。换言之，自同构
$J_0:\mathbb R^{2n}\to\mathbb R^{2n}$ 由下列 $2n\times2n$ 矩阵表示：
\[
J_0:=
\begin{pmatrix}
0&-1&\cdots&0&0\\
1&0&\cdots&0&0\\
\vdots&\vdots&\ddots&\vdots&\vdots\\
0&0&\cdots&0&-1\\
0&0&\cdots&1&0
\end{pmatrix}.
\]
每个\BookRecordIndex{T0207}凯勒流形都给出另一个例子（但并非每个\BookRecordIndex{T0358}辛流形都容许凯勒结构）。

\paragraph{\BookRecordIndex{T0003}$J$-全纯曲线。}
\zhindex{j-quanqvxian}{\texorpdfstring{$J$-全纯曲线}{J-全纯曲线}}
一条 \BookRecordIndex{T0003}$J$-全纯曲线，是从黎曼曲面 $(\Sigma,j)$ 到殆复流形 $(M,J)$ 的
$(j,J)$-全纯映射
\[
  u:\Sigma\to M.
\]
通常取 $(\Sigma,j)$ 为黎曼球面 $S^2=\mathbb C\cup\{\infty\}$。
\footnote{黎曼曲面是复维数为 $1$ 的复流形。}
依照复几何的术语，把 $2$-球面看成复射影直线 $\mathbb{CP}^1$ 往往很方便。

若 $u$ 是嵌入（即单射浸入），则其像 $C$ 是 $M$ 的一个 $2$ 维子流形，且切空间
$T_xC$ 在 $J$ 下不变。因此每个切空间都是 $TM$ 中的一条复直线。再由\BookRecordIndex{T0372}驯服条件，
$\omega$ 在每条这样的直线上的限制都是正形式，所以 $C$ 是 $M$ 的辛子流形。
\footnote{若 $\omega$ 在 $M$ 的子流形 $X$ 上的限制非退化，就称 $X$ 为辛子流形。}
反过来，若 $C\subset M$ 是 $2$ 维辛子流形，则存在一个被 $\omega$ \BookRecordIndex{T0372}驯服的\BookRecordIndex{T0088}殆复结构
$J$，使 $TC$ 在 $J$ 下不变。（先在 $TC$ 上定义 $J$，再把它延拓到所有
$x\in C$ 处的切空间 $T_xM$，最后把这个截面延拓到 $M$ 的其余部分。）由于 $J$
在 $TC$ 上的限制可积，可知 $C$ 是一条 \BookRecordIndex{T0003}$J$-全纯曲线的像。因此，嵌入的
\BookRecordIndex{T0003}$J$-全纯曲线本质上与 $M$ 的 $2$ 维辛子流形相同。这一论证本质地用到了 $C$ 的
$2$ 维性；从高维流形到 $(M,J)$ 的 $J$-全纯映射不可能形成一种良好的理论。（第 12.7 节
提到的\BookRecordIndex{T0356}辛涡旋方程并不与此说法矛盾，因为那些方程带有额外结构。）

请注意，按照这个定义，曲线 $u$ 总是带参数的。这一点应与复几何中的情形对照：在复几何
中，人们常把曲线定义为若干全纯多项式的公共零点集。这种做法在不可积的殆复世界里没有
意义，因为当 $J$ 不可积时，通常不存在全纯函数 $(M,J)\to\mathbb C$。
\footnote{不过，Donaldson 最近发展出一套强有力的渐近全纯函数
$(M,J,\omega)\to\mathbb C$ 理论；参见 Donaldson \srcite{84,85} 和
Auroux \srcite{26}。}
